\chapter{Implementation of the Margin}
\label{chap:implementation}

% TODO: Describe how \cref{chap:margin} was turned into working
% code. Suggested structure:
%  - Software stack / tools (language, libraries, snake robot
%    simulator, solvers used for the wrench-space computation).
%  - Algorithm(s): pseudocode for computing the margin from a
%    given set of contacts (use the algorithm/algpseudocode
%    environment below).
%  - Numerical considerations: tolerances, degenerate contact
%    configurations, computational cost / complexity in practice.
%  - Validation of the implementation against known/simple cases
%    before using it in simulation (e.g. hand-computable examples).
%  - Integration into the path planner, if implemented.
%
% STRUCTURE (restructured 2026-09-30): Software and Tools, Algorithm.
% Moved OUT of this chapter:
%  - The LP2 formulation, lambda ordering, and the explanation of d and
%    of the normalization e^T lambda <= n_c -> theory.tex, Sec. LP2
%    (eq:lp2). Only the solver-specific conversion is kept here.
% Removed (empty, not in the new outline): "Numerical Considerations",
% "Verification of the Implementation" -- their intended content is in
% the TODO list above.

\section{Software and Tools}
\label{sec:software}

% TODO: language, libraries, solver (e.g. scipy.optimize.linprog), simulator.

\section{Algorithm}
\label{sec:algorithm}

% [First paragraph kept from implementation.tex, section "Algorithm";
%  the rest of the LP2 derivation that followed it was moved to
%  theory.tex, Sec. LP2.]
To evaluate the force-closure quality of a given contact configuration, the
linear program LP2~\eqref{eq:lp2} from \cite[Sec.~28.4.2]{prattichizzo2008grasping}
is employed. The optimization determines a force-closure margin $d$ by
finding a feasible contact-force distribution that maximizes the minimum
margin to the boundaries of the contact constraints.

\subsection{Conversion to Standard LP Form}

The optimization is implemented using a linear-programming solver in the
standard form
\begin{equation}
\begin{aligned}
    \underset{\boldsymbol{x}}{\operatorname{minimize}}
    \quad & \boldsymbol{c}^{T}\boldsymbol{x}\\
    \text{subject to}
    \quad & A_{\mathrm{ub}}\boldsymbol{x}
        \leq \boldsymbol{b}_{\mathrm{ub}},\\
    & A_{\mathrm{eq}}\boldsymbol{x}
        = \boldsymbol{b}_{\mathrm{eq}},
\end{aligned}
\end{equation}
with appropriate variable bounds.

Since the optimization variable contains both the force intensities and the
force-closure margin, it is defined as
\begin{equation}
    \boldsymbol{x}
    =
    \begin{bmatrix}
        \boldsymbol{\lambda}\\
        d
    \end{bmatrix}.
\end{equation}

Most linear-programming solvers, including the solver used in this work,
formulate the objective as a minimization. Therefore, maximizing $d$ is
implemented by minimizing $-d$. The objective vector is consequently
\begin{equation}
    \boldsymbol{c}
    =
    \begin{bmatrix}
        \boldsymbol{0}\\
        -1
    \end{bmatrix},
\end{equation}
such that
\begin{equation}
    \boldsymbol{c}^{T}\boldsymbol{x}=-d.
\end{equation}

The equilibrium constraint is written as
\begin{equation}
    A_{\mathrm{eq}}
    =
    \begin{bmatrix}
        G & \boldsymbol{0}
    \end{bmatrix},
    \qquad
    \boldsymbol{b}_{\mathrm{eq}}
    =
    \boldsymbol{0}.
\end{equation}

The constraint
\begin{equation}
    F\boldsymbol{\lambda}
    -\boldsymbol{1}d
    \geq \boldsymbol{0}
\end{equation}
must be converted to the solver's required ``less-than-or-equal-to''
form. Multiplication by $-1$ gives
\begin{equation}
    -F\boldsymbol{\lambda}
    +\boldsymbol{1}d
    \leq \boldsymbol{0}.
\end{equation}
Consequently, the corresponding inequality matrix is
\begin{equation}
    A_{\mathrm{ub}}
    =
    \begin{bmatrix}
        -F & \boldsymbol{1}\\
        \boldsymbol{e}^{T} & 0
    \end{bmatrix},
\end{equation}
with
\begin{equation}
    \boldsymbol{b}_{\mathrm{ub}}
    =
    \begin{bmatrix}
        \boldsymbol{0}\\
        n_c
    \end{bmatrix}.
\end{equation}

The optimization variables are subject to the bounds
\begin{equation}
    -\infty < \lambda_i < \infty,
    \qquad
    d\geq0.
\end{equation}

The resulting LP is therefore
\begin{equation}
\boxed{
\begin{aligned}
    \underset{\boldsymbol{x}}{\operatorname{minimize}}
    \quad&
    \begin{bmatrix}
        \boldsymbol{0}^{T} & -1
    \end{bmatrix}
    \boldsymbol{x}\\
    \text{subject to}
    \quad&
    \begin{bmatrix}
        -F & \boldsymbol{1}\\
        \boldsymbol{e}^{T} & 0
    \end{bmatrix}
    \boldsymbol{x}
    \leq
    \begin{bmatrix}
        \boldsymbol{0}\\
        n_c
    \end{bmatrix},\\
    &
    \begin{bmatrix}
        G & \boldsymbol{0}
    \end{bmatrix}
    \boldsymbol{x}
    =
    \boldsymbol{0},\\
    &d\geq0.
\end{aligned}
}
\end{equation}

The optimal value of the original maximization problem is recovered from
the LP solver by changing the sign of the objective value,
\begin{equation}
    d^{*}=-\boldsymbol{c}^{T}\boldsymbol{x}^{*}.
\end{equation}
A positive value of $d^{*}$ indicates that a contact-force distribution
satisfying the equilibrium and contact constraints exists with a positive
margin. The magnitude of $d^{*}$ consequently provides a measure of the
force-closure margin for the given contact configuration.

\begin{algorithm}[H]
\caption{Force closure margin computation (placeholder)}
\begin{algorithmic}[1]
\Require Contact points $\{p_i\}$, contact normals $\{n_i\}$, friction coefficients $\{\mu_i\}$
\Ensure Force closure margin $m$
\State Build wrench cone for each contact $i$
\State Assemble total wrench space $W$
\State Compute margin $m \gets$ \Call{ComputeMargin}{$W$}
\State \Return $m$
\end{algorithmic}
\end{algorithm}
